%======================================================================
\section{Computations across the hierarchy}
\label{sec:iterative}

In this section we place examples of computations at every level of the
hierarchy
\[
\text{stable}\Rightarrow\NIP_3\Rightarrow\NIP_2\Rightarrow\NIP_1\Rightarrow\NIP,
\]
and show that each implication is strict.
For each level we give two examples.
The first is one of the static examples of Section~\ref{sec:examples},
turned into a CCS by the register construction of
Definition~\ref{def:register}; its deep computations are one of the five
compacta of Section~\ref{sec:five-compacta}.
The second is a classical iterative algorithm, a single step map applied over
and over.
For iterative algorithms, the hierarchy coincides with a hierarchy studied in
topological dynamics by Glasner and Megrelishvili~\cite{glasner2022tame}, and
we begin with this observation.
Table~\ref{tab:hierarchy} at the end of the section summarizes the examples.

\subsection{Iterative algorithms and tame dynamics}
A \emph{cascade CCS} is a compositional computation structure
$(L,\mathcal{P},\Gamma)$ in which $L$ is a compact metrizable space,
$\mathcal{P}$ is a countable family of bounded continuous predicates that
separates the points of~$L$, and $\Gamma=\{\gamma^n:n\in\mathbb{N}\}$ is the
semigroup of iterates of a single continuous map $\gamma\colon L\to L$.
Since $\mathcal{P}$ separates points and $L$ is compact, the map
$v\mapsto\tp(v)$ is a homeomorphism of $L$ onto the space of types, so $L$ is
covered by a single shard, every type is realized, and the Extendibility Axiom
holds with $\tilde\gamma=\gamma$.
The deep computations of a cascade CCS form the pointwise closure
\[
E(L,\gamma)=\overline{\{\gamma^n:n\in\mathbb{N}\}}\subseteq L^L ,
\]
which is the \emph{enveloping semigroup} of the dynamical system $(L,\gamma)$
introduced by Ellis.
A dynamical system is called \emph{tame} if its enveloping semigroup consists of
Baire class~$1$ maps~\cite{Kohler1995,Glasner2006}.
Theorem~\ref{thm:nip-baire1def} therefore reads as follows.

\begin{prop}\label{prop:cascade-tame}
A cascade CCS is $\NIP$ if and only if the dynamical system $(L,\gamma)$ is
tame.
In that case it is $\NIP_1$, $\NIP_2$ or $\NIP_3$ if and only if $E(L,\gamma)$ is,
respectively, first countable, hereditarily separable or metrizable.
\end{prop}

For metrizable systems, Glasner and Megrelishvili call these three classes
\emph{tame$_1$}, \emph{tame$_2$} and \emph{hereditarily non-sensitive}, and
prove that they are strictly nested~\cite[Proposition~0.6]{glasner2022tame}.
The dynamical point of view supplies a criterion for wildness.
By a theorem of Kerr and Li~\cite{KerrLi2007}, a tame system given by a
homeomorphism has topological entropy zero.
An algorithm that produces information at a positive exponential rate, such as
a pseudorandom number generator, is wild by design; a numerical method that
settles down is tame.


\subsection{Stable and metrizable}
\label{sec:ex-NIP3}
\emph{A stable example: rotation.}
Let $L=\mathbb{T}=\mathbb{R}/\mathbb{Z}$, with the predicates $\cos2\pi x$ and
$\sin2\pi x$, and let $\gamma(x)=x+\alpha$ with $\alpha$ irrational.
A pointwise limit of rotations is a rotation, and the rotations $n\alpha$ are
dense, so the deep computations are all the rotations
$x\mapsto x+c$, $c\in\mathbb{T}$.
They are continuous, so the rotation CCS is stable, and its space of deep
computations is the circle group, a metrizable compactum.

\emph{$\NIP_3$ but not stable: square roots.}
The square-root CCS of Section~\ref{exa:square-roots-compute} has the deep
computations $\Gamma\cup\{\mathfrak f\}$, a countable compactum, hence
metrizable; and $\mathfrak f$ is discontinuous at $0$ and $\infty$.
So it is $\NIP_3$ but not stable.

\emph{$\NIP_3$ but not stable: the power method.}
The power method computes a dominant eigenvector of a matrix by repeated
multiplication and normalization.
Let $A=\operatorname{diag}(\lambda_1,\dots,\lambda_d)$ with
$\lambda_1>\lambda_2>\cdots>\lambda_d>0$, let $L=\RP^{d-1}$, and let
$\mathcal{P}$ consist of the predicates
$p_{ij}([v])=v_iv_j/\|v\|^2$, the entries of the orthogonal projection onto the
line $[v]$.
The step of the algorithm is $\gamma([v])=[Av]$.
For $[v]\in\RP^{d-1}$, let $i(v)$ be the least index with $v_i\ne0$, and put
$\pi([v])=[e_{i(v)}]$.

\begin{prop}\label{prop:power}
The deep computations of the power method are
$E(L,\gamma)=\{\gamma^n:n\in\mathbb{N}\}\cup\{\pi\}$.
Hence the power method is $\NIP_3$.
Its only nontrivial deep computation $\pi$ is of Baire class~$1$ and
discontinuous, so the power method is not stable.
\end{prop}

\begin{proof}
For $i=i(v)$, we have $\lambda_i^{-n}A^nv=v_ie_i+\sum_{j>i}(\lambda_j/\lambda_i)^nv_je_j\to v_ie_i$,
so $\gamma^n\to\pi$ pointwise.
Every cluster point of $\{\gamma^n\}$ is therefore either some $\gamma^m$ or
$\pi$, and $E(L,\gamma)$ is a countable compact Hausdorff space, hence
metrizable.
The map $\pi$ is a pointwise limit of continuous maps, and it is
discontinuous at every $[v]$ with $i(v)>1$, since $[v+te_1]\to[v]$ as $t\to0$
while $\pi([v+te_1])=[e_1]\ne\pi([v])$.
\end{proof}

The deep computation $\pi$, \emph{the dominant eigenvector of the input}, is
not computable in the sense of Section~\ref{sec:CCS}, since it does not extend
continuously to the space of types; yet it is the limit of a sequence of
computations, and the closure in which it lives is as simple as possible.


\subsection{Hereditarily separable but not metrizable}
\label{sec:ex-NIP2}
\emph{Threshold classifiers.}
By Lemma~\ref{lem:register}, the finite-precision threshold classifiers of
Section~\ref{sec:Cantor-classifiers} form a register CCS whose deep
computations are the ultra-classifiers.
Their space is a copy of $S(2^{\mathbb{N}})$
(Section~\ref{sec:infin-prec-class}), so this CCS is $\NIP_2$ but not $\NIP_3$.
The classifiers are $\NIP$ because the sets they define are initial segments of
the lexicographic order, and no infinite family of initial segments is
independent.

\emph{Drawing a line.}
Let $0<\alpha<1$ be irrational.
To draw the line $y=\alpha x+\rho$ on a screen, Bresenham's
algorithm~\cite{Bresenham1965} moves one pixel to the right at each step and
decides whether to move one pixel up.
The $n$-th decision is
$b_n(\rho)=\lfloor(n+1)\alpha+\rho\rfloor-\lfloor n\alpha+\rho\rfloor
=\Brack{\{n\alpha+\rho\}\in[1-\alpha,1)}$, where $\{\cdot\}$ denotes the
fractional part, and the state of the algorithm is the error term $\rho$
modulo~$1$.
The step $\rho\mapsto\rho+\alpha$ is a rotation of the circle
$\mathbb{T}=\mathbb{R}/\mathbb{Z}$, but the decision bit is not a continuous
function of~$\rho$.
The standard remedy~\cite{MorseHedlund1940,Lothaire2002} is to take as states
the decision sequences themselves.

Let $D=\{-k\alpha:k\in\mathbb{Z}\}\subseteq\mathbb{T}$, a countable dense set
containing both endpoints $0$ and $1-\alpha$ of the decision interval.
The \emph{split circle} $\mathbb{T}_D$ is obtained from $\mathbb{T}$ by replacing
each $t\in D$ by two points $t^-<t^+$, with the circular order topology; it is
a compact metrizable space, homeomorphic to the Cantor set.
The rotation by $\alpha$ preserves $D$ and lifts to a homeomorphism $R$ of
$\mathbb{T}_D$ that preserves the labels $\pm$.
The decision interval becomes clopen in $\mathbb{T}_D$, and the coding map
$\mathbb{T}_D\to2^{\mathbb{Z}}$, $y\mapsto(\Brack{R^ny\in[1-\alpha,1)})_{n\in\mathbb{Z}}$,
is a homeomorphism onto the \emph{Sturmian shift} $X_\alpha$ that conjugates $R$
to the shift~$\sigma$.
The \emph{Bresenham CCS} has states $L=X_\alpha$, the coordinate predicates, and
the semigroup $\Gamma=\{\sigma^n:n\in\mathbb{N}\}$.

For $c\in\mathbb{T}$, let $R_c^+$ (respectively $R_c^-$) be the map of
$\mathbb{T}_D$ that sends a point over $y_0$ to $(y_0+c)^+$ (respectively
$(y_0+c)^-$) when $y_0+c\in D$, and to $y_0+c$ otherwise.
The \emph{double arrow circle} $\mathbb{D}$ is $\mathbb{T}\times\{-,+\}$ with the
lexicographic circular order and its order topology; it is compact, first
countable, hereditarily separable and not metrizable~\cite[3.10.C]{Engelking:1989}.

\begin{thm}\label{thm:Bresenham}
The Bresenham CCS is $\NIP$.
Its deep computations are
\[
E(X_\alpha,\sigma)=\{\sigma^n:n\in\mathbb{N}\}\cup\{R_c^+,R_c^-:c\in\mathbb{T}\},
\]
where each $\sigma^n$ is isolated and $c^\pm\mapsto R_c^\pm$ is a homeomorphism
of $\mathbb{D}$ onto the remaining part.
In particular, the Bresenham CCS is $\NIP_2$ but not $\NIP_3$.
\end{thm}

\begin{proof}
For $\NIP$, note that the $j$-th coordinate of $\sigma^n y$ is $1$ exactly when
$R^{n+j}y$ lies in the clopen arc $[1-\alpha,1)$ of $\mathbb{T}_D$, that is, when
$y$ lies in an arc of $\mathbb{T}_D$.
A family of arcs of a circularly ordered set has Vapnik--Chervonenkis dimension
at most~$3$, so it has no infinite independent subfamily, and the predicates
$P_j\circ\sigma^n$ are $\{0,1\}$-valued.

For the deep computations, let $\gamma$ be the limit of $\sigma^{n_i}$ along an
ultrafilter, and let $c=\lim n_i\alpha$ in $\mathbb{T}$.
If $n_i$ is eventually constant, then $\gamma$ is a power of $\sigma$.
Otherwise, since the points $n\alpha$ are distinct, the ultrafilter decides
whether $n_i\alpha$ approaches $c$ from the right or from the left.
In the first case, for every point $y$ over $y_0$ the points $R^{n_i}y$ over
$y_0+n_i\alpha$ approach $y_0+c$ strictly from the right, and they converge
in $\mathbb{T}_D$ to $(y_0+c)^+$ if $y_0+c\in D$ and to $y_0+c$ otherwise; so
$\gamma=R_c^+$.
The second case gives $R_c^-$.
Conversely, every $c$ is approached from both sides by the points $n\alpha$
with $n\ge N$, for every $N$, so every $R_c^\pm$ is a deep computation, and
$R_c^+\ne R_c^-$, since they differ at the points over $d-c$ with $d\in D$.
The same argument shows that $c_i^{\varepsilon_i}\to c^+$ in $\mathbb{D}$ implies
$R_{c_i}^{\varepsilon_i}\to R_c^+$ pointwise, and symmetrically for $c^-$; so
$c^\pm\mapsto R_c^\pm$ is a continuous bijection from the compact space
$\mathbb{D}$ onto a Hausdorff space, hence a homeomorphism.
Finally, fix $d\in D$ and consider the two points
$(d-n\alpha)^\pm$ of $\mathbb{T}_D$, which $\sigma^n$ sends to $d^-$ and $d^+$.
Every map $R_c^\pm$ sends them to the same point, and so does every pointwise
limit of $\sigma^{m_i}$ with $m_i\to\infty$, by the argument above.
A net of such maps, or of maps $\sigma^m$ with $m\ne n$ taken from a bounded set of exponents, cannot converge to $\sigma^n$, so $\sigma^n$ is isolated.
Thus $E(X_\alpha,\sigma)$ is the union of a countable set and a copy of
$\mathbb{D}$; it is hereditarily separable and not metrizable.
\end{proof}

The two deep computations $R_c^+$ and $R_c^-$ are the two ways of drawing the
line infinitely far away, where the accumulated error has reached $c$: they
agree except at the pixels where the line passes exactly through a corner,
and there they apply the two opposite rounding conventions.
Figure~\ref{fig:bresenham} shows the two conventions.

\begin{figure}[htbp]
  \centering
  \includegraphics[width=\textwidth]{bresenham.pdf}
  \caption{Left: the pixels drawn by Bresenham’s algorithm for $\alpha=(\sqrt5-1)/2$, with the decision bits $b_n$.
    Right: where the line passes exactly through a pixel corner, the deep computations $R_c^+$ and $R_c^-$ break the tie in opposite directions; they agree on every other pixel.}
  \label{fig:bresenham}
\end{figure}
The split Cantor set of the threshold classifiers of
Section~\ref{sec:Cantor-classifiers} is the same phenomenon in a static form.


\subsection{First countable but not hereditarily separable}
\label{sec:ex-NIP1}
\emph{Three-way comparison.}
For $t\in2^{<\mathbb{N}}$, let $h_t\colon2^{\mathbb{N}}\to\{0,\frac12,1\}$ take the
value $0$ below the interval $[t^\frown0^\infty,t^\frown1^\infty]$, the value
$\frac12$ on it, and $1$ above it.
This is one step of a binary search with three outcomes, as in a search tree:
the input lies below, inside or above the dyadic interval of~$t$.
Each $h_t$ is continuous, and its superlevel sets are final segments of the
lexicographic order, so the family $\{h_t\}$ is $\NIP$.
The map $h_t\mapsto(v_t,f^+_{t^\frown0^\infty})$ extends to a homeomorphism of
the closure of $\{h_t\}$ onto $\hat D(S(2^{\mathbb{N}}))$
\cite[4.3.7]{argyros2008rosenthal}.
By Lemma~\ref{lem:register}, the register CCS of the three-way comparisons is
$\NIP_1$ but not $\NIP_2$.

\emph{Drawing a line in space.}
Let $\alpha$ and $\beta$ be irrational numbers in $(0,1)$ such that $1,\alpha,\beta$
are linearly independent over $\mathbb{Q}$.
To draw a line in $\mathbb{Z}^3$ with direction $(1,\alpha,\beta)$, the
three-dimensional Bresenham algorithm runs the planar algorithm in two
coordinates at once.
Its CCS has states $L=X_\alpha\times X_\beta$, the coordinate predicates of
both factors, and the semigroup generated by $\sigma\times\sigma$.

\begin{thm}\label{thm:Bresenham3}
The three-dimensional Bresenham CCS is $\NIP_1$ but not $\NIP_2$.
\end{thm}

\begin{proof}
Each predicate is a coordinate of one factor, so the CCS is $\NIP$ by
Theorem~\ref{thm:Bresenham}.
The deep computations form a closed subspace of
$E(X_\alpha,\sigma)\times E(X_\beta,\sigma)$, which is first countable as a
product of two first countable spaces; so the CCS is $\NIP_1$.
By Kronecker's theorem, the points $(n\alpha,n\beta)$ with $n\ge N$ are dense in
$\mathbb{T}^2$ for every $N$, so every point $(c,c')$ is approached by them from
inside each of the four open quadrants at $(c,c')$; arguing as in the proof of
Theorem~\ref{thm:Bresenham}, every pair $(R_c^\varepsilon,R_{c'}^{\varepsilon'})$
is a deep computation.
Consider the pairs $(R_c^+,R_{-c}^+)$, for $c\in\mathbb{T}$.
In $\mathbb{D}\times\mathbb{D}$, a neighborhood of $(c^+,(-c)^+)$ is
$[c^+,(c+\delta)^-)\times[(-c)^+,(-c+\delta)^-)$, and it contains
$(s^+,(-s)^+)$ only if $c\le s<c+\delta$ and $-c\le -s<-c+\delta$, that is,
only if $s=c$.
So these pairs form an uncountable discrete subspace, and the space of deep
computations is not hereditarily separable.
\end{proof}

The uncountable discrete set consists of the deep computations that break
ties upward in both coordinates at points where the two accumulated errors are
opposite.
Each of the two planar algorithms is $\NIP_2$, and drawing in two directions at
once leaves the class: like the Sorgenfrey plane, a product of two
hereditarily separable spaces need not be hereditarily separable.
Compare \cite[Proposition~4.1]{glasner2022tame}, where a product of two
tame$_2$ systems is shown not to be tame$_2$.


\subsection{NIP but not first countable}
\label{sec:ex-NIP}
\emph{Prefix tests.}
By Lemma~\ref{lem:register}, the prefix tests of Section~\ref{sec:prefix-test}
form a register CCS whose deep computations are the ultra-prefix tests, a copy
of $\hat A(2^{\mathbb{N}})$ (Section~\ref{sec:extended-alexandroff}).
Any two prefix tests are either nested or disjoint, so no infinite family of
them is independent, and the CCS is $\NIP$ but not $\NIP_1$.

\emph{Tolerance tests.}
Floating-point code does not test real numbers for equality; it tests whether
$|x-a|$ is below a tolerance.
For a rational $q\in[0,1]$ and $n\ge1$, let $h_{q,n}(x)=\max(0,1-n|x-q|)$, the
degree to which $x$ equals $q$ at precision $1/n$.
The \emph{tolerance CCS} is the register CCS (Definition~\ref{def:register}) of
the family $\{h_{q,n}\}$ on $X=[0,1]$: the \emph{tolerance test} $\gamma_{q,n}$
writes $h_{q,n}(x)$ into the register.

\begin{prop}\label{prop:tolerance}
The tolerance CCS is $\NIP$ but not $\NIP_1$.
Its deep computations include the exact equality tests
$(x,y)\mapsto(x,\Brack{x=a})$ for every real $a\in[0,1]$, and the test
$(x,y)\mapsto(x,0)$ that never succeeds.
\end{prop}

\begin{proof}
By Lemma~\ref{lem:register}, it suffices to consider the functions $h_{q,n}$.
For $\NIP$, the sets $\{h_{q,n}\ge b\}$ are intervals, and an independent
sequence of the $h_{q,n}$ would give an independent sequence of intervals, which
does not exist.
Given $a$, choose rationals $q_k\to a$ and integers $n_k\to\infty$ with
$n_k|q_k-a|\to0$; then $h_{q_k,n_k}(a)\to1$ and $h_{q_k,n_k}(x)\to0$ for
$x\ne a$, so $\Brack{x=a}$ is a deep computation.
Exact equality tests at distinct points $a_k$ converge pointwise to~$0$.
The pointwise topology on $\{\Brack{\,\cdot=a}:a\in[0,1]\}\cup\{0\}$ is that of the
Alexandroff compactification of a discrete set of size continuum, which is not
first countable at~$0$; hence neither is the space of deep computations.
\end{proof}

Exact equality is thus the deep computation of tolerance testing (Figure~\ref{fig:tolerance}).

\begin{figure}[htbp]
  \centering
  \includegraphics[width=0.78\textwidth]{tolerance.pdf}
  \caption{Tolerance tests $h_{q,n}$ with $q\to a$ and $n\to\infty$, chosen so that $n|q-a|\to0$, converge pointwise to the exact equality test $\Brack{x=a}$, which is $1$ at $a$ and $0$ elsewhere.}
  \label{fig:tolerance}
\end{figure}
The failure of first countability has a concrete meaning: a neighborhood of
the test that never succeeds is specified by finitely many inputs at which the
test must fail, and no countable family of such specifications suffices,
because exact tests can succeed at uncountably many inputs.


\subsection{Beyond NIP: Newton's method on the Riemann sphere}
Newton's method on the Riemann sphere is chaotic on its Julia set, and the
chaos is visible to the predicates of the CCS.
For a polynomial $p$ of degree $m\ge2$ with simple roots, the Newton map
$N_p$ is a rational map of degree~$m$, whose topological entropy is
$\log m>0$~\cite{Lyubich1983}, and positive entropy is incompatible with
tameness for invertible systems~\cite{KerrLi2007}.
The Newton map is not invertible, so we do not claim the general statement here;
for square roots the independence can be exhibited by hand.

\begin{prop}\label{prop:Newton-wild}
Let $p(z)=z^2-1$, and let $N(z)=\frac12\bigl(z+\frac1z\bigr)$ be its Newton map on
$\hatC$.
The functions $P_2\circ N^{3k}$, $k\in\mathbb{N}$, form an independent sequence:
for all finite disjoint sets $E,F\subseteq\mathbb{N}$ there is $z\in\hatC$ with
$P_2(N^{3k}z)\ge\frac12$ for $k\in E$ and $P_2(N^{3k}z)\le-\frac12$ for $k\in F$.
In particular, the Newton CCS of $p$ is not $\NIP$.
\end{prop}

\begin{proof}
The Möbius transformation $\phi(z)=(z-1)/(z+1)$ conjugates $N$ to the squaring
map: $\phi(N(z))=\phi(z)^2$.
It maps the imaginary axis, together with $\infty$, onto the unit circle, with
$\phi^{-1}(e^{i\theta})=i\cot(\theta/2)$.
Hence, for $z=i\cot(\theta/2)$, we have $N^n(z)=i\cot(2^{n}\theta/2)$ and
\[
P_2\bigl(N^n(z)\bigr)=\frac{2\cot(2^{n-1}\theta)}{\cot^2(2^{n-1}\theta)+1}=\sin(2^n\theta).
\]
Write $\theta=2\pi\sum_{j\ge1}x_j2^{-j}$ with bits $x_j$.
The fractional part of $2^{3k}\theta/2\pi$ has binary expansion
$0.x_{3k+1}x_{3k+2}x_{3k+3}\cdots$.
If $(x_{3k+1},x_{3k+2},x_{3k+3})=(0,1,0)$, this fractional part lies in
$[\frac14,\frac38]$, so $\sin(2^{3k}\theta)\ge\sin\frac{3\pi}{4}>\frac12$; if the block
is $(1,1,0)$, it lies in $[\frac34,\frac78]$, so $\sin(2^{3k}\theta)<-\frac12$.
Choosing the block $(0,1,0)$ for $k\in E$ and $(1,1,0)$ for $k\in F$ gives the
required~$z$.
\end{proof}

The same algorithm is tame on the real line.
The Julia set of $N$ is the imaginary axis together with $\infty$, which meets
$\RP^1$ only at~$\infty$, and on $\RP^1$ the iterates converge to the deep
computation of Example~\ref{exa:square-roots-compute}.
Whether Newton's method is tame depends on the space of states, not on the
formula for the step.


\subsection{Summary}
Table~\ref{tab:hierarchy} collects the examples of this section.

\begin{table}[ht]
\centering
\small
\renewcommand{\arraystretch}{1.25}
\begin{tabular}{>{\raggedright\arraybackslash}p{2.3cm}>{\raggedright\arraybackslash}p{2.7cm}>{\raggedright\arraybackslash}p{2.6cm}>{\raggedright\arraybackslash}p{3.0cm}}
\hline
Level & Static example & Iterative algorithm & Deep computations contain\\
\hline
stable & --- & rotation & the circle group\\
$\NIP_3$, not stable & square roots & power method & a countable compactum\\
$\NIP_2$, not $\NIP_3$ & threshold classifiers & Bresenham, plane & $S(2^{\mathbb{N}})$; a double arrow circle\\
$\NIP_1$, not $\NIP_2$ & three-way comparison & Bresenham, space & $\hat D(S(2^{\mathbb{N}}))$; $\mathbb{D}\times\mathbb{D}$\\
$\NIP$, not $\NIP_1$ & prefix tests & --- & $\hat A(2^{\mathbb{N}})$\\
& tolerance tests & & $A(2^{\mathbb{N}})$\\
not $\NIP$ & --- & Newton on $\hatC$ & ---\\
\hline
\end{tabular}
\vspace{6pt}
\caption{Computations at each level of the hierarchy.}
\label{tab:hierarchy}
\end{table}
